% 对应题号 M2-116，本题条件、变量及几何关系逐项核对。
% 修订：启用中文断行；核对题设及图文对应。
\documentclass[tikz,border=14pt]{standalone}
\usepackage{fontspec,xeCJK}
\usepackage{amsmath,amssymb}
\usepackage{xcolor}
\usepackage{tikz}
\usetikzlibrary{arrows.meta,calc,patterns,positioning,bending}
\IfFontExistsTF{Heiti SC}{\IfFontExistsTF{Heiti SC}{\setmainfont{Heiti SC}}{\setmainfont{FandolHei-Regular}}\IfFontExistsTF{Heiti SC}{\setsansfont{Heiti SC}}{\setsansfont{TeX Gyre Heros}}}{\IfFontExistsTF{Songti SC}{\setmainfont{Songti SC}}{\setmainfont{FandolSong-Regular}}\IfFontExistsTF{Songti SC}{\setsansfont{Songti SC}}{\setsansfont{FandolSong-Regular}}}
\definecolor{bluefield}{RGB}{24,76,255}
\definecolor{redfield}{RGB}{232,49,35}
\definecolor{orangefield}{RGB}{236,132,18}
\definecolor{greenfield}{RGB}{18,132,86}
\definecolor{purplefield}{RGB}{112,72,190}
\definecolor{linegray}{RGB}{60,68,78}
\definecolor{softblue}{RGB}{238,243,255}
\definecolor{softgreen}{RGB}{238,249,244}
\definecolor{softorange}{RGB}{255,247,235}
\tikzset{
  arrow/.style={-{Latex[length=3.1mm,width=2.1mm]},line width=1.1pt},
  thinarr/.style={-{Latex[length=2.5mm,width=1.7mm]},line width=.85pt},
  axis/.style={-{Latex[length=2.5mm,width=1.7mm]},line width=.8pt,draw=linegray},
  guide/.style={draw=linegray!55,dashed,line width=.65pt},
  box/.style={draw=linegray,rounded corners=3pt,fill=white,align=center,inner sep=6pt,line width=.75pt},
  bluebox/.style={box,draw=bluefield,fill=softblue},
  greenbox/.style={box,draw=greenfield,fill=softgreen},
  orangebox/.style={box,draw=orangefield,fill=softorange}
}
\IfFontExistsTF{Heiti SC}{\setCJKmainfont[AutoFakeBold=1.5]{Heiti SC}}{\setCJKmainfont[AutoFakeBold=1.5]{FandolHei-Regular}}
\IfFontExistsTF{Heiti SC}{\setCJKsansfont{Heiti SC}}{\setCJKsansfont{FandolHei-Regular}}
\begin{document}

\begin{tikzpicture}[x=1cm,y=1cm]
\path[use as bounding box] (-10,-4.6) rectangle (10,4.6);
\fill[white] (-10,-4.6) rectangle (10,4.6);
\node[font=\Large\bfseries,text width=18.6cm,align=center] at (0,4.05) {矩形三边接地、上边等势：角点跳变不影响内部唯一解};
\begin{scope}[shift={(-4.8,-.15)}]
\fill[softblue] (0,0) rectangle (4.4,3); \draw[linegray,line width=1.3pt] (0,0)--(4.4,0)--(4.4,3); \draw[linegray,line width=1.3pt] (0,0)--(0,3);
\draw[redfield,line width=2pt] (.08,3)--(4.32,3); \node[redfield,above] at (2.2,3) {$V_0$};
\node[below] at (2.2,0) {$0$}; \node[left] at (0,1.5) {$0$}; \node[right] at (4.4,1.5) {$0$};
\fill[orangefield] (0,3) circle(.09); \fill[orangefield] (4.4,3) circle(.09); \node[orangefield] at (2.2,-.7) {两上角为边界跳变点};
\end{scope}
\begin{scope}[shift={(4.65,-.1)}]
\node[bluebox,text width=6.9cm] at (0,1.65) {$\displaystyle \varphi=\sum_{n=1}^{\infty}A_n\sin\frac{n\pi x}{a}\,\frac{\sinh(n\pi y/a)}{\sinh(n\pi b/a)}$};
\node[greenbox,text width=6.9cm] at (0,0) {$A_n=\dfrac{2V_0}{n\pi}[1-(-1)^n]$，故仅奇数模态出现。};
\node[orangebox,text width=6.4cm,align=center] at (0,-1.7) {角点的单点赋值不改变内部调和解；\\正弦级数在端点取零，域内趋近角点无唯一极限。};
\end{scope}
\end{tikzpicture}
\end{document}
